\section{Closed-Loop Multitrophic Ecosystems: Mycoponics, Bioweathering, and Substrate Regeneration}
\label{sec:closed_loop}

\subsection{Saprotrophic Fungal Bioconversion of Inedible Biomass}
\label{subsec:mycoponics}

Primary candidate space crops---including dwarf wheat (\textit{Triticum aestivum}), potato (\textit{Solanum tuberosum}), and soybean (\textit{Glycine max})---exhibit harvest indices between $0.40$ and $0.60$ under optimized controlled environment agriculture\cite{wheeler2017agriculture,ming1989lunar,wheeler1996crop}. Consequently, $40\%$ to $60\%$ of total produced crop biomass consists of inedible lignocellulosic residues, primarily fibrous stalks, roots, and seed chaff\cite{ming1989lunar,wheeler1996crop,gitelson1995biological}. If left untreated, this biomass immobilizes carbon, nitrogen, and phosphorus, breaking closed-loop nutrient cycling within the habitat\cite{wheeler2017agriculture,ming1989lunar,liu2020bioregenerative}.

Integrating saprotrophic fungi, such as oyster mushrooms (\textit{Pleurotus ostreatus}) and Reishi (\textit{Ganoderma lucidum}), into the life-support loop resolves this elemental bottleneck\cite{ming1989lunar,liu2020bioregenerative,eaglin2026clinostat}. Fungi secrete robust extracellular enzyme suites---including laccases, manganese peroxidases, and endo-$\beta$-glucanases---that break down recalcitrant lignocellulose matrices into simple soluble carbohydrates\cite{ming1989lunar,krause2020microbial}. This bioconversion produces protein- and vitamin-rich edible fruiting bodies, directly increasing the habitat's effective harvest index and dietary diversity while respiring metabolic $\mathrm{CO}_2$ back to the crop canopy\cite{ming1989lunar,liu2020bioregenerative}.

However, fungal cultivation cannot be accomplished by simply incorporating straw into raw lunar regolith\cite{ming1989lunar,eaglin2026clinostat}. Unweathered mineral shards physically slice expanding fungal hyphae, and exposed mineral coordination sites adsorb and inactivate extracellular enzymes\cite{turci2015mechanistic,fackrell2024mechanical}. Sintered porous ceramics provide an ideal inert substrate platform for mycoponics\cite{ming1989lunar}. Perforated ceramic trays and vertical columns can be loaded with pasteurized crop residues and inoculated with fungal spawn\cite{ming1989lunar}. The smooth, passivated ceramic pore surfaces protect delicate hyphal networks, maintain optimal moisture retention, and preserve enzyme catalytic activity, enabling reliable fungal bioconversion in closed loops\cite{ming1989lunar,eaglin2026clinostat}.

\subsection{Microbial Bioweathering for Closed-Loop Mineral Replenishment}
\label{subsec:bioweathering}

Although sintered ceramics remain stable in hydroponic solution, continuous crop harvesting removes essential macronutrients ($\mathrm{K}^+, \mathrm{Ca}^{2+}, \mathrm{Mg}^{2+}, \mathrm{PO}_4^{3-}$) and trace metals from the liquid loop\cite{ming1989lunar,atkin2024bioremediation,cockell2018biorock}. To prevent indefinite resupply of terrestrial fertilizer salts, these nutrients can be extracted in situ from raw regolith via microbial bioleaching\cite{ming1989lunar,cockell2018biorock,santomartino2024ecomine}.

Rather than exposing plant roots directly to abrasive dust, bioweathering is conducted in dedicated external bioreactors\cite{ming1989lunar,cockell2018biorock}. The feasibility of microbially mediated mineral extraction in space was established by the European Space Agency's BioRock and BioAsteroid flight experiments aboard the International Space Station, which demonstrated enhanced element recovery from basaltic materials across microgravity, Martian gravity, and terrestrial gravity regimes\cite{cockell2018biorock,cockell2020microbially}. Microorganisms such as \textit{Sphingomonas desiccabilis}, \textit{Bacillus subtilis}, and \textit{Penicillium simplicissimum} mobilize structural minerals via three distinct pathways\cite{krause2020microbial,cockell2018biorock,santomartino2024ecomine}:
\begin{enumerate}
    \item \textbf{Acidolysis:} Secretion of low-molecular-weight organic acids (primarily citric, oxalic, and gluconic acids) lowers solution $\mathrm{pH}$, protonating bridging oxygen atoms and accelerating the dissolution of silicate glasses, apatite, and ilmenite\cite{krause2020microbial,cockell2018biorock,santomartino2024ecomine}.
    \item \textbf{Complexolysis:} Organic acid anions chelate polyvalent cations ($\mathrm{Fe}^{3+}, \mathrm{Al}^{3+}, \mathrm{Ca}^{2+}, \mathrm{Mg}^{2+}$), forming stable organometallic complexes that maintain cations in solution and prevent secondary mineral precipitation\cite{fateri2019solar}.
    \item \textbf{Siderophore-Mediated Iron Chelation:} Bacteria secrete specialized high-affinity iron-chelating ligands (siderophores) that extract ferric iron directly from mineral lattices, even under neutral-to-alkaline conditions\cite{cockell2018biorock,raafat2022biomining}.
\end{enumerate}

Following leaching, the leachate undergoes ultrafiltration and ion-exchange separation to remove bacterial biomass, aluminum, and heavy metal precipitates before introduction into the primary hydroponic loop\cite{ming1989lunar}. This architecture decouples mineral extraction from crop production, obtaining essential plant nutrients locally while shielding crops from the physical and geochemical hazards of raw lunar dust\cite{paul2022plants,ming1989lunar}.

\subsection{Thermal Substrate Regeneration and Materials Upcycling}
\label{subsec:thermal_regeneration}

Over continuous cultivation cycles, hydroponic substrates accumulate senescent root fragments, organic root exudates, and microbial biofilms\cite{ming1989lunar}. In terrestrial greenhouse operations, organic fouling clogs pore networks and harbors phytopathogens, necessitating the disposal and replacement of substrates such as rockwool\cite{ming1989lunar,wamelink2019crop}. Sintered ceramics, by contrast, possess the thermal durability required for cyclic regeneration\cite{ming1989lunar,barta1994thermal}.

Following harvest, root-bound ceramic modules undergo pyrolytic oxidation at $500^\circ\mathrm{C}$ in a controlled oxygen atmosphere\cite{ming1989lunar,barta1994thermal}. At this temperature, the aluminosilicate ceramic substrate remains dimensionally and microstructurally stable, whereas all accumulated organic root biomass, extracellular polysaccharides, and biofilms are pyrolyzed and oxidized\cite{ming1989lunar,barta1994thermal}:
\begin{equation}
    \mathrm{C}_x\mathrm{H}_y\mathrm{O}_z (\text{biomass}) + \left(x + \frac{y}{4} - \frac{z}{2}\right)\mathrm{O}_2 \xrightarrow{500^\circ\mathrm{C}} x\,\mathrm{CO}_2 + \frac{y}{2}\,\mathrm{H}_2\mathrm{O} + \text{Mineral Ash}
    \label{eq:pyrolytic_bakeout}
\end{equation}

The volatilized carbon dioxide and water vapor are directed into the habitat ECLSS loop; moisture is recovered via condensing heat exchangers, while $\mathrm{CO}_2$ is returned to the crop canopy for photosynthesis\cite{wheeler2017agriculture,ming1989lunar,barta1994thermal}. The residual mineral ash, containing non-volatile potassium, phosphorus, and trace nutrients, is washed out and redissolved into the liquid fertilizer inventory\cite{ming1989lunar,schoonen2020hydroxyl}. The ceramic matrix emerges sterile, decarbonized, and with unclogged pore channels, ready for immediate replanting without structural degradation\cite{ming1989lunar}.

When ceramic modules eventually reach their mechanical fatigue life, they can be crushed and upcycled as gravel beds for habitat drainage systems, re-ground into raw feedstock for secondary sintering, or utilized as rigid structural scaffolds for mycelium-bound biocomposite construction materials\cite{ming1989lunar}.
